% Appendix K -- Compute and Reproducibility (body after the \ifarxiv ... \fi sentence, which stays in main.tex).
% Sources: main/reports/FINAL_A.md 8.2-8.4, main/LOG.md, rss_ext/LOG.md, tmlr_ext/LOG.md and RESULTS_tmlr.md 4 (keys k.*,
%   kind 'report'); DESIGN_tmlr.md 0.2. Anonymous: no names, server, repository or home paths (PLAN 4.2 item 6, 5.3).
\textbf{Compute.} The main study (E0--E7) ran on two RTX A5000 GPUs of a local server and occupied them for about \nv{k.gpuh.main} GPU-hours: \nv{k.gpuh.td} hours of one GPU shared by 16 MPS processes for TD-MPC2 and \nv{k.gpuh.dw} GPU-hours for DINO-WM, including time slowed by sharing. The theory simulations took \nv{k.cpuh.e4} CPU-hours. The extension E8--E10 with E8c ran on rented A40, RTX 4090 and RTX 5090 GPUs for \$\nv{k.cost.rss}, and E11--E17 with the post hoc runs on rented A40 and RTX 4090 GPUs for \$\nv{k.cost.tmlr} (\$\nv{k.cost.tmlr.main} for the planned runs and \$\nv{k.cost.tmlr.posthoc} for D4--D6), including setup and idle time; X6 (D7) cost a further \$\nv{k.cost.x6}. Every job of E11--E17 completed; jobs hit by run incidents were re-run (Appendix~\ref{app:plan}).

\textbf{What the supplementary material contains.} The code of every stage as it ran; the design documents of all versions, two documents of them redacted (the main design and the initial experiment plan; Appendix~\ref{app:plan}); the coordinator and worker logs; the stage reports; and the result files the scripts read and the per-stage analysis outputs, including episode-level tables for the TD-MPC2 stages E9--E17. The raw job records of the cloud machines (about 210 MB) are not included. It is organized as the main study (designs, code, results per stage, reports), the extension E8--E10, and E11--E17 with one folder per group of runs (TD-MPC2 5M, 317M, DINO-WM, integer kernel), each holding its operational plan fixed before the main run, its code, results, log and report; a separate folder holds the scripts that recomputed the verdicts and the re-analyses R1--R5. Checkpoints and datasets are not included: they are the public releases of TD-MPC2 and DINO-WM, and their hashes are recorded. Upstream commits and software versions are listed in Appendix~\ref{app:details}.

\textbf{Recomputing every number.} Every result number in this paper is a macro defined in \texttt{numbers.tex}, which \texttt{build/make\_numbers.py} writes from the result files; each definition records the file and field it comes from. \texttt{build/make\_tables.py} writes every numeric table from the same files. \texttt{build/verify\_numbers.py} re-runs the computation and checks that \texttt{numbers.tex} is current, that every macro used in the text is defined, that statements whose wording depends on the data (``in all cells'', ``none harmful'', ``indistinguishable from zero'') agree with it, and that every number still typed in the text is a registered design constant or label. A few values exist only in the stage reports and logs (the compute and cost figures above and the number of simulated planner runs); \texttt{numbers.tex} marks them, and \texttt{verify\_numbers.py} lists them as not verified against data. Run from \texttt{paper/tmlr}, the commands \texttt{python3 build/make\_numbers.py}, \texttt{python3 build/make\_tables.py} and \texttt{python3 build/verify\_numbers.py --strict} regenerate and check all numbers from the supplementary material alone, without a GPU; these scripts use only the Python standard library (with a UTF-8 locale), and the figures additionally need matplotlib and PyMuPDF. The supplement includes an anonymized copy of the manuscript source, which \texttt{verify\_numbers.py} reads, and a README that maps every table and figure to its result files. The episode-level tables of E9--E17 on TD-MPC2 allow their intervals to be recomputed (paired bootstrap, $10^4$ resamples, NumPy \texttt{RandomState(0)}); for E1--E7 and for DINO-WM the supplement holds the per-cell and per-job summaries, whose episode records are part of the raw job records. The verdicts of E11, E14, K2 and the DINO-WM stages were recomputed from the raw job records by scripts independent of the analysis code. These scripts and their outputs are included, but the raw records they read are not.

\textbf{Re-running the experiments.} Seeds are listed in Appendix~\ref{app:details}. Results are bit-identical only on the same GPU type and software stack: on a second machine of the same type, the E11 cells of \detnoise and of the unquantized model were reproduced bit for bit, the first jobs of every new DINO-WM machine matched those of the first one bit for bit, and on another GPU type the planner outputs differ while the episode outcomes agree in most episodes (X17, Appendix~\ref{app:new}).
